\documentclass[11pt,a4paper,utf8x]{report}
\usepackage[T1]{fontenc} 
\usepackage[utf8]{inputenc} 
\usepackage{lmodern}
\usepackage [frenchb]{babel}
% Pour pouvoir utiliser 
\usepackage{ucs}

\usepackage{textcomp}
\usepackage{graphicx}
\usepackage{keystroke}
\usepackage{amssymb}
\usepackage{amsmath}
%\renewcommand{\thesection}{\arabic{section}} % numérotation des sections	
\usepackage[cc]{titlepic} %rajouter le logo dans la page de garde
\usepackage{url} % Pour avoir de belles url
\usepackage {geometry}
\usepackage[linktocpage]{hyperref}
% Pour pouvoir faire plusieurs colonnes
\usepackage {multicol}
%Pour faire plusieurs lignes
\usepackage{multirow}
\usepackage{slashbox}
% Pour mettre du code source
\usepackage {listings}
% Pour pouvoir passer en paysage
\usepackage{lscape}	

% Pour pouvoir faire plusieurs colonnes
\usepackage {multicol}

% POur crééer un index
\usepackage{makeidx}

%Pour taper des algorithme en pseudo code
%\usepackage{algorithm,algorithmic}

%insertion de pdf
\usepackage{pdfpages}

\hypersetup{
backref=true,
%permet d'ajouter des liens dans...
pagebackref=true,%...les bibliographies
hyperindex=true, %ajoute des liens dans les index.
colorlinks=true, %colorise les liens
breaklinks=true, %permet le retour à la ligne dans les liens trop longs
urlcolor= blue, %couleur des hyperliens
citecolor=	cyan,
bookmarks=true, %créé des signets pour Acrobat
bookmarksopen=true,
%si les signets Acrobat sont créés,
%les afficher complètement.
pdftitle={Initiation à la Recherche}, %informations apparaissant dans
pdfauthor={MARGUERITE Alain\\ RINCE Romain},
%dans les informations du document
pdfsubject={Doc}
%sous Acrobat.
}

\makeindex


%%%% debut macro pour enlever le nom chapitre %%%%
\makeatletter
\def\@makechapterhead#1{%
  \vspace*{50\p@}%
  {\parindent \z@ \raggedright \normalfont
    \interlinepenalty\@M
    \ifnum \c@secnumdepth >\m@ne
        \Huge\bfseries \thechapter\quad
    \fi
    \Huge \bfseries #1\par\nobreak
    \vskip 40\p@
  }}

\def\@makeschapterhead#1{%
  \vspace*{50\p@}%
  {\parindent \z@ \raggedright
    \normalfont
    \interlinepenalty\@M
    \Huge \bfseries  #1\par\nobreak
    \vskip 40\p@
  }}
\makeatother
%%%% fin macro %%%%

%Couverture 

%Couverture 

\widowpenalty=10000
\clubpenalty=10000
 
\title{ {\huge \textsc{Components and Architectures}} \\Master ALMA 2\up{eme} année \\\vspace{3cm} \emph{TP Composant : Implémentation HADL}\vspace{3cm}}

\titlepic{ \includegraphics[scale=0.80]{img/logouniv}}


\author{A.\textsc{Marguerite}\\ R.\textsc{Rincé}\vspace{3cm}}
\date{Université de Nantes \\ 2 rue de la Houssinière, BP92208, F-44322 Nantes cedex 03, FRANCE}
%\vspace{20cm}


\newcommand{\cf}{\emph{cf}.}

%\vspace{9cm}
\date{Université de Nantes \\ 2 rue de la Houssinière, BP92208, F-44322 Nantes cedex 03, FRANCE}
%\vspace{20cm}
\begin{document}

\maketitle
\renewcommand{\labelitemi}{$\bullet$} 

\clearpage

\tableofcontents
\clearpage

%\chapter*{Remerciements}
\phantomsection
\addcontentsline{toc}{section}{Remerciements}
 TODO TODO TODO TODO TODO TODO TODO TODO TODO TODO TODO TODO TODO TODO TODO TODO TODO TODO TODO TODO TODO TODO

 TODO TODO TODO TODO TODO TODO TODO TODO TODO TODO TODO TODO TODO TODO TODO TODO TODO TODO TODO TODO

\chapter{Introduction}\label{chap:Intro}

\section{Problématiques}
Les problèmes d'inter-opérabilité et de réutilisabilité font partie des problématiques majeures des systèmes informatiques. La programmation objet ne parvient pas à répondre à ces problématiques à cause notamment de ses propriétés de granularité et de fort couplage. Les problèmes de cohésions (pas d'explications des besoins requis) sont un frein à l'utilisation de la programmation objet à grande échelle. 

Parmi les solutions à ces problématiques les architectures à composants sont des alternatives efficaces. Un composant, entité principale, correspond par définition à une ou plusieurs fonctionnalités. Il est donc possible de le concevoir et de le pré-tester de manière indépendante. La coopération entres composant est alors plus aisée. De plus un composant possède un couplage lâche. Cette propriété garantissant l'inter-opératibilté est permise par l'utilisation d'interfaces entre composants. 

\section{Présentation du sujet}
Le module \textsc{Architectures and Components} nous propose une étude de la programmation par composants. L'objectif de ce Projet/TP  est de réaliser un Home Architecture Description Language, et d'implémenter une architecture classique \emph{Clients-Serveur}. La démarche proposée pour atteindre cet objectif est composée de trois étape : 


\begin{enumerate}
\item 
Réalisation du M2 :\hfill \\
Il s'agit d'un méta-modèle de l'architecture en composants.  
\item
Réalisation du M1 :\hfill \\
À partir du méta-modèle précédement réalisé, nous pouvons proposer une instance de ce dernier et obtenir un modèle. Nous choisirons ici de modéliser une application classique \emph{Client-Serveurs}.  
\item
Réalisation du M0 :\hfill \\
Dernière étape, elle consiste à l'implémentation du M1 et de coder concrétement qui y sont modélisées.
\end{enumerate}

Dans ce rapport nous détaillerons chaque étape la modélisation proposée puis l'implémentation.

\chapter{Implémentation du M1}
Dans ce chapitre nous définirons une application \textit{Client-Serveur} en implémentant le méta-model proposé au chapitre \ref{chap:M2}.

\section{Démarche modélisation M1}
L'énoncé du projet/TP nous propose un serveur avec les fonctionnalités suivantes :

\begin{itemize}
\item \verb+ConnexionManager+ :   Gestion des connexions. 
\item  \verb+Database+ :  Gestion de la base de données.
\item \verb+SecurityManager+ :  Gestion d'un système de sécurité.
\end{itemize}

Un composant étant par définition la spécification d'une fonctionnalités, chacune des fonctionnalités précedement listées est un \verb+Composant+. Chaque service de chaque composant a son port requis relié au rôle fournis  du connecteur ( et inversement pour les ports fournis et rôles requis).

Les liaisons est ces différents composant et le client sont explicitées dans la figure \ref{fig:desSer}.
\begin{figure}[htb]
  \centering
  \includegraphics[scale=0.32]{img/DescribServeur}
  \caption{Description du Serveur}
  \label{fig:desSer}.
\end{figure}


\subsection{ConnexionManager}
Ce composant possède les services suivants : 
\begin{itemize}
\item 
  \verb+ServiceFExternalSocket+ : Fournit une connexion au serveur.
\item 
  \verb+ServicerDDBQuery+ :  requière une réponse de la base de données (requête SQL).
\item 
  \verb+ServiceRSecurityManager+ :  requière une validation de l'authentification
\end{itemize}


\subsection{Database}
Ce composant possède les services suivants : 
\begin{itemize}
\item 
  \verb+ServiceFSecurityManagement+ :  fournit une réponse à la requête de sécurité.
\item 
  \verb+ServicerFExecuteSQL+ :  fournit une réponse à une requête du client (par le ConnexionManager).
\end{itemize}

\subsection{SecurityManager}
Ce composant possède les services suivants : 
\begin{itemize}
\item 
  \verb+ServiceFSecurityAuth+ : fournit une validation de la demande de connexion.
\item 
  \verb+ServicerSecurityManager+ :  requière une réponde de sécurité de la base de données (requête SQL).
\end{itemize}

\section{Diagramme M1}
\pagestyle{empty}
%\newgeometry{a3paper}
\begin{figure}[htb]
  \centering
  \includegraphics[scale=0.20]{img/M1}
  \caption{Model (M1)}
  \label{fig:M1}
\end{figure}

\chapter{Implémentation du M2}\label{chap:M2}
\section{Détails de la modélisation}
Construisons un méta-modèle basé sur l'entité \composant :
\begin{itemize}
\item 
  Un \composant{} possède des \proprietes{}. Ces dernières peuvent être fonctionelles (\verb+PropFonc+) ou non fonctionelles (\verb+PropNonFonc+).
\item 
  Il possède deux interfaces composants :  une pour les \services{} requis, une pour les \verb+Services+ fournis.
\item
  Une interface est aussi constituée de \ports{} (fournis ou requis). 
\item
  Chaque port est associé à un et un seul \service, mais un \service{} possède plusieurs \ports.  
\end{itemize}


\paragraph{}
Afin d'effectuer la communication entres composants, il est nécessaire d'introduire la notion de connecteur :

\begin{itemize}
\item 
  Un \connecteur{} a deux \verb+InterfaceConnecteur+.
\item
  Chaque interface connecteur possède deux rôles : un rôle requis \roler{} et un rôle fournios\rolef{}.
\item
  Chaque rôle d'une interface est relié par un lien \attachement{} à son port correspondant (\textit{c.a.d }: un rôle requis à un port requis et un rôle fournis à un port fournis).  
\item
  Les deux interfaces d'un connecteur sont reliés par deux \glu{}, chaque \glu{} reliant un rôle requis au rôle fournis opposé.
\end{itemize}

\paragraph{}

La \configuration{} est un composant composé de \verb+composants+ : 

\begin{itemize}
\item 
  Elle possède deux \verb+InterfaceConfig+. Ces dernières hériant d'\verb+InterfaceComposant+.
\item
  Il est possible de relier un \portr{} de l'\verb+InterfaceConfig+ à un \portr{} d'un composant de la configuration (et réciproquement pour le \portf) grâce à un lien \verb+binding+.
\item
  Un connecteur peut être aussi une \configuration.
\end{itemize}
%% \begin{description}
%% \item[G] \hfill \\
%%   Mélanges de classes de dépendances.
%% \item[Passage à l'échelle] \hfill \\
%%   Au bout d'un grand nombre de classe l'évolution et l'extension se voient limitées.
%% \end{description}
\ssection{Diagramme M2}
\pagestyle{empty}
\begin{figure}[htb]
  %\centering
  \includegraphics[scale=0.36]{img/M2}
  \caption{Méta-Model (M2)}
  \label{fig:M2}
\end{figure}


\listoffigures
% Pour finir l'interligne de 1,5
%\bibliographystyle{alpha}
%\bibliography{biblio.bib}


\end{document}
